\chapter{Geriatrics \& Gerontology}

\medterm{Age-Related Macular Degeneration (AMD)}-- Progressive degenerative disease of the central retina (macula) causing bilateral, irreversible loss of central vision. Types: dry (atrophic, 90%) with drusen accumulation; wet (neovascular, 10%) with choroidal neovascularization. Risk factors: age >50, smoking, family history, Caucasians. Symptoms: central scotoma, metamorphopsia, decreased color vision. Diagnosis: ophthalmoscopy, Amsler grid, OCT, fluorescein angiography. Treatment: dry--AREDS2 supplements (vitamins C, E, zinc, copper, lutein, zeaxanthin); wet--anti-VEGF injections (ranibizumab, aflibercept, bevacizumab).

\medterm{Barthel Index}-- Standardized measure of functional independence in activities of daily living (ADL). Scores 0-100 based on 10 items: feeding, bathing, grooming, dressing, bowel control, bladder control, toilet use, transfers, mobility, stair climbing. <40: severe dependency; 41-60: moderate dependency; >60: mild dependency. Used for care planning, discharge decisions, outcome measurement.

\medterm{Benign Prostatic Hyperplasia (BPH)}-- Non-cancerous enlargement of prostate gland causing lower urinary tract symptoms (LUTS) in aging men. Prevalence: ~50% at age 60, ~90% at age 85. Symptoms: hesitancy, weak stream, urgency, frequency, nocturia, incomplete emptying. Diagnosis: DRE, PSA, uroflowmetry, post-void residual. Treatment: watchful waiting, alpha-blockers (tamsulosin), 5-alpha-reductase inhibitors (finasteride), combination therapy, surgical (TURP, laser).

\medterm{Bioethics (Geriatric)}-- Ethical considerations in elderly care: autonomy vs. paternalism, informed consent capacity, end-of-life decisions, resource allocation, elder abuse. Principles: respect for autonomy, beneficence, non-maleficence, justice. Advance directives: living will, durable power of attorney for healthcare. Surrogate decision-makers when capacity lacking.

\medterm{Bone Density}-- Measurement of bone mineral content (BMC) and density (BMD). Assessed by DEXA scan (gold standard). T-score: compares to young adult reference. Normal: >-1.0; osteopenia: -1.0 to -2.5; osteoporosis: <-2.5. Fracture risk increases with decreasing BMD. Screening: women >65, men >70, or younger with risk factors.

\medterm{Catheter-Associated UTI (CAUTI)}-- Urinary tract infection related to indwelling urinary catheter. Most common healthcare-associated infection. Risk: prolonged catheterization, female sex, diabetes, immunosuppression. Prevention: avoid unnecessary catheters, aseptic insertion, closed drainage, daily hygiene, early removal. Treatment: antibiotics based on culture, catheter removal or exchange.

\medterm{Cognitive Impairment}-- Decline in one or more cognitive domains: memory, attention, language, visuospatial, executive function. Types: subjective cognitive decline, mild cognitive impairment (MCI), dementia. MCI: functional independence preserved, impairment on testing. Dementia: functional impairment. Assessment: MMSE, MoCA, detailed neuropsychological testing.

\medterm{Comorbidity}-- Presence of two or more chronic conditions in same individual. Common in elderly: hypertension, diabetes, arthritis, COPD, depression. Increases complexity of care, polypharmacy risk, functional decline. Comprehensive assessment needed for each condition and interactions.

\medterm{Confusion (Delirium)}-- Acute onset fluctuating disturbance in attention and awareness. Usually reversible. Precipitants: infection, medications, metabolic derangement, pain, constipation, urinary retention. Differentiate from dementia (chronic, progressive). Treatment: identify and treat underlying cause, environmental modification, avoid anticholinergics.

\medterm{Deconditioning}-- Physiological decline from reduced activity or immobilization. Effects: muscle atrophy, weakness, decreased cardiorespiratory fitness, balance impairment, increased fall risk. Preventable with early mobilization, exercise. Reversible with rehabilitation.

\medterm{Dementia}-- Chronic progressive cognitive decline interfering with daily function. Types: Alzheimer disease (60-70%), vascular dementia (10-20%), Lewy body dementia (10-15%), frontotemporal dementia (5-10%), mixed. Symptoms: memory loss, aphasia, apraxia, agnosia, executive dysfunction, behavioral changes. Diagnosis: clinical, cognitive testing, imaging, lab workup. Treatment: symptomatic (cholinesterase inhibitors, memantine), manage behavioral symptoms.

\medterm{Dentistry (Geriatric)}-- Oral health challenges in elderly: xerostomia, root caries, periodontal disease, denture problems, oral cancer, medication side effects. Assessment: complete oral exam, cancer screening, denture fit, caries risk. Prevention: fluoride, oral hygiene, regular dental visits. Adaptations: handholds, magnification, positioning.

\medterm{Depression (Geriatric)}-- Mood disorder common in elderly, often underdiagnosed. Symptoms: sad mood, anhedonia, sleep/appetite changes, fatigue, guilt, concentration difficulty, suicidal ideation. May present as cognitive impairment ("pseudodementia"). Risk factors: chronic illness, social isolation, bereavement, medications. Treatment: psychotherapy, antidepressants (SSRIs first-line), social support. Screen: Geriatric Depression Scale (GDS-15).

\medterm{Disease-Modifying Antirheumatic Drug (DMARD)}-- See rheumatology.

\medterm{Drift (Sit-to-Stand)}-- Functional test measuring lower extremity strength and balance. Timed Up and Go (TUG): stand, walk 3 meters, turn, return, sit. Normal: <10 seconds; fall risk: >12 seconds. Three-step squat test, chair rise test alternative.

\medterm{Dysphagia (Geriatric)}-- Difficulty swallowing, common in elderly. Causes: neurological (stroke, Parkinson's), structural (cancer, stricture), muscular (presbyesophagus), medications. Symptoms: coughing with eating, choking, pocketing food, weight loss, aspiration. Assessment: clinical exam, videofluoroscopy, endoscopy. Management: diet modification, swallowing therapy, positioning, feeding tube if severe.

\medterm{Fall}-- Unintended change in position to ground. Major cause of morbidity/mortality in elderly. Risk factors: gait/balance impairment, muscle weakness, visual impairment, polypharmacy, environmental hazards, orthostatic hypotension. Prevention: exercise, medication review, vision correction, home modification, hip protectors. Assessment: gait, balance, strength, vision, medications, environment.

\medterm{Frailty}-- Clinical syndrome of decreased reserve and resilience to stressors. Components: unintentional weight loss, exhaustion, low physical activity, slow gait speed, weak grip strength. Identification: Fried phenotype, frailty index. Consequences: adverse outcomes (falls, hospitalization, mortality), polypharmacy, functional decline. Intervention: exercise, nutrition, medication review, comprehensive geriatric assessment.

\medterm{Functional Status}-- Ability to perform activities of daily living (ADLs) and instrumental ADLs (IADLs). ADLs: bathing, dressing, toileting, transferring, continence, feeding. IADLs: shopping, cooking, managing medications, finances, transportation, housekeeping. Decline indicates need for assistance.

\medterm{Geriatric Assessment}-- Comprehensive multidimensional evaluation of elderly person. Domains: medical, functional, cognitive, psychological, social, environmental, pharmacological. Performed by interdisciplinary team. Guides care planning, identifies problems, predicts outcomes.

\medterm{Glycemic Control}-- Maintenance of blood glucose within target range. In elderly, goals less stringent to avoid hypoglycemia. HbA1c target: <7.5-8.0% depending on health status. Monitoring: self-monitoring, continuous glucose monitoring. Medications: metformin first-line, avoid sulfonylureas if hypoglycemia risk.

\medterm{Health Promotion}-- Activities maintaining/improving health in elderly. Components: vaccination, screening, exercise, nutrition, cognitive stimulation, social engagement, fall prevention. Tailored to individual capacity and preferences.

\medterm{Hip Fracture}-- Fracture of proximal femur, usually femoral neck or intertrochanteric region. Mechanism: low-energy fall in osteoporotic bone. Symptoms: pain, inability to bear weight, shortened/externally rotated leg. Diagnosis: X-ray, MRI if X-ray negative. Treatment: surgical fixation (internal fixation, arthroplasty). Complications: DVT, PE, pneumonia, pressure ulcers, mortality (20-30% at 1 year).

\medterm{Hospitalization (Geriatric)}-- Inpatient stay for elderly patients. High risk for adverse outcomes: delirium, functional decline, institutionalization. Precautions: fall prevention, early mobilization, medication reconciliation, delirium prevention, nutrition support. Geriatrician consultation beneficial.

\medterm{Incontinence (Urinary)}-- Involuntary urine leakage. Types: stress (cough/sneeze), urge (overactive bladder), overflow (retention), functional (immobility/cognition), mixed. Assessment: history, bladder diary, urinalysis, post-void residual. Treatment: behavioral (bladder training, pelvic floor), pharmacological (anticholinergics, beta-3 agonists), surgical (sling, artificial sphincter).

\medterm{Incontinence (Fecal)}-- Involuntary fecal leakage. Causes: diarrhea, constipation with overflow, sphincter damage, neurological disease, medications. Assessment: history, examination, anorectal manometry. Treatment: bowel regimen, sphincter exercises, biofeedback, surgery if indicated.

\medterm{Infection (Geriatric)}-- Infectious diseases in elderly often present atypically: fever may be absent, confusion may be sole symptom. Common infections: UTI, pneumonia, skin/soft tissue, bloodstream. Prevention: vaccination, hand hygiene, aspiration precautions. Treatment: earlier presentation, broader differential, drug interactions, renal dosing.

\medterm{Instrumental ADL}-- See functional status.

\medterm{Isolation (Social)}-- Lack of social contacts and interactions. Common in elderly: bereavement, mobility limitation, sensory impairment, retirement. Consequences: depression, cognitive decline, functional impairment, mortality. Interventions: social programs, transportation, technology, volunteer visits.

\medterm{Medication Review}-- Systematic examination of all medications for appropriateness. Components: indication, efficacy, safety, adherence, interactions. Tools: Beers Criteria, STOPP/START. Goals: deprescribing, optimizing therapy, reducing adverse effects.

\medterm{Mortality}-- Death rate in population. Elderly mortality: higher rates, common causes: cardiovascular disease, cancer, dementia, respiratory disease, fall-related injuries. Life expectancy varies by sex, race, socioeconomic status.

\medterm{Neurocognitive Disorder}-- See dementia.

\medterm{Nutrition (Geriatric)}-- Nutritional assessment and intervention in elderly. Risks: malnutrition (undernutrition, obesity), micronutrient deficiencies (vitamin D, B12, calcium), dehydration. Assessment: weight, BMI, dietary intake, lab values. Interventions: supplementation, meal delivery, appetite stimulants, feeding assistance.

\medterm{Osteoporosis}-- Systemic skeletal disease characterized by low bone mass and microarchitectural deterioration. Results in increased fracture risk. Diagnosis: DEXA T-score <-2.5. Risk factors: age, female sex, Caucasian/Asian, low body weight, smoking, alcohol, glucocorticoids, family history. Prevention: calcium, vitamin D, exercise, fall prevention. Treatment: bisphosphonates, denosumab, teriparatide, calcitonin.

\medterm{Polypharmacy}-- Use of multiple medications (usually 5+). Common in elderly due to multiple chronic conditions. Risks: adverse drug reactions, interactions, non-adherence, functional decline. Strategies: deprescribing, medication review, simplification, patient education.

\medterm{Pressure Ulcer}-- See pressure injury.

\medterm{Pressure Injury}-- Localized damage to skin and underlying tissue, usually over bony prominence. Stages: 1 (non-blanchable erythema), 2 (partial thickness loss), 3 (full thickness loss), 4 (full thickness with exposed bone/tendon), unstageable (obscured), deep tissue injury. Prevention: repositioning, support surfaces, nutrition, skin care.

\medterm{Prosthetics}-- Artificial replacement of missing body part. Elderly use: lower limb (transtibial, transfemoral), upper limb. Considerations: mobility level, cognition, vision, comorbidities, support system. Rehabilitation: physical therapy, strength training, balance.

\medterm{Palliative Care}-- Specialized medical care for serious illness focusing on symptom relief and quality of life. Not limited to end-of-life. Can be provided alongside curative treatment. Components: pain management, symptom control, psychosocial support, spiritual care, advance care planning. Referral early in disease course.

\medterm{Sarcopenia}-- Age-related loss of muscle mass and function. Prevalence increases with age. Risk factors: reduced protein intake, decreased physical activity, chronic disease. Assessment: grip strength, gait speed, muscle mass measurement. Treatment: protein supplementation, resistance exercise, vitamin D, hormonal therapy if deficient.

\medterm{Social Determinants of Health}-- Conditions in which people are born, grow, live, work, age. Include: income, education, employment, housing, social support, physical environment. Major influence on health equity. Elderly vulnerable: fixed income, isolation, functional limitations.

\medterm{Survival}-- Length of life after diagnosis or event. In elderly, consider comorbidities, functional status, preferences. Life expectancy estimates based on age, sex, health status. Shared decision-making essential.

\medterm{Swallowing}-- See dysphagia.

\medterm{Urinary Retention}-- Inability to empty bladder completely. Acute: painful, emergency. Chronic: painless, overflow incontinence. Causes: BPH, medications (anticholinergics), neurological disease, obstruction. Diagnosis: post-void residual >100-150 mL. Treatment: catheterization, alpha-blockers, surgery for BPH.

\medterm{Vision Impairment}-- Reduced visual acuity or field. Common causes: cataract, AMD, glaucoma, diabetic retinopathy, stroke. Screening: annual eye exam after 65. Interventions: cataract surgery, glaucoma treatment, low vision aids. Impact: functional limitation, fall risk, isolation.

\medterm{Weight Loss}-- Unintentional weight loss >5% in 6-12 months. Common in elderly, associated with morbidity/mortality. Causes: malignancy, depression, dementia, dental problems, medications, chronic disease, socioeconomic factors. Assessment: history, examination, targeted testing. Treatment: address underlying cause, nutritional support, appetite stimulants.

\medterm{Xerostomia}-- Dry mouth due to decreased saliva. Common in elderly from medications (anticholinergics, diuretics, antidepressants), Sjogren syndrome, radiation. Symptoms: difficulty swallowing, speaking, tasting, dental caries. Treatment: hydration, sugarless gum/candy, saliva substitutes, pilocarpine.

\medterm{Yield (Geriatric)}-- In geriatric context, yield refers to diagnostic yield of tests or intervention effectiveness. Screening yield varies by age, risk factors. Balance benefits vs. harms in elderly.

\medterm{Zollinger-Ellison Syndrome}-- See gastroenterology.

\medterm{Advance Directive}-- Legal document expressing healthcare wishes. Components: living will (treatment preferences), durable power of attorney for healthcare (surrogate decision-maker). Should be reviewed periodically, discussed with surrogate and healthcare team.

\medterm{Catastrophic Aging}-- Rapid functional decline after stressor (illness, surgery, fall). Elderly have decreased reserve, less ability to recover. Prevention: comprehensive assessment, optimization, early intervention.

\medterm{Delirium}-- See confusion.

\medterm{Elder Abuse}-- Intentional or neglectful act causing harm to older person. Types: physical, emotional, sexual, financial, neglect. Risk factors: social isolation, caregiver stress, cognitive impairment, dependency. Mandatory reporting required. Assessment: direct question, observation, collateral history.

\medterm{Functional Assessment}-- Evaluation of ability to perform daily activities. Tools: Barthel Index, Lawton IADL Scale, Katz Index. Identifies needs for assistance, predicts outcomes, guides care planning.

\medterm{Geriatric Syndrome}-- Multifactorial condition common in elderly, not attributable to single disease. Examples: falls, incontinence, frailty, delirium, immobility, pressure ulcers. Require comprehensive assessment and multidisciplinary management.

\medterm{Hospice}-- Palliative care for terminal illness with prognosis <6 months. Focus on comfort, quality of life, support for patient and family. Provided at home, hospice facility, or hospital. Interdisciplinary team: physician, nurse, social worker, chaplain, volunteer.

\medterm{Informed Consent (Geriatric)}-- Process of obtaining voluntary authorization from elderly patient with decision-making capacity. Capacity assessment: understanding, appreciation, reasoning, communication of choice. Impaired capacity may require surrogate decision-maker.

\medterm{Long-Term Care}-- Assistance with ADLs/IADLs for extended period. Settings: home, assisted living, nursing home. Funding: Medicaid, long-term care insurance, out-of-pocket. Planning essential for elderly and families.

\medterm{Memory Loss}-- See cognitive impairment/dementia.

\medterm{Nutritional Screening}-- Rapid assessment for malnutrition risk. Tools: MST, MNA, MUST. Screening all elderly patients recommended. Identifies need for comprehensive assessment.

\medterm{Orthostatic Hypotension}-- Blood pressure drop on standing: systolic >20 mmHg or diastolic >10 mmHg within 3 minutes. Causes: volume depletion, medications, autonomic dysfunction, prolonged bed rest. Symptoms: dizziness, lightheadedness, syncope, falls. Treatment: non-pharmacological (compression stockings, gradual position changes), pharmacological (fludrocortisone, midodrine).

\medterm{Polypharmacy}-- See polypharmacy.

\medterm{Restraint}-- Device or medication limiting movement. Physical restraints: vests, wrist restraints, lap traps. Chemical restraints: antipsychotics, benzodiazepines for behavior control. Should be last resort, time-limited, with monitoring. Alternatives: addressing underlying causes, environmental modification.

\medterm{Syncope}-- Transient loss of consciousness due to cerebral hypoperfusion. Causes: vasovagal, orthostatic, cardiac (arrhythmia, structural), neurologic. Assessment: history, examination, ECG, orthostatic BP. Risk stratification: San Francisco Syncope Rule, OESIL score.

\medterm{Well-Older Person}-- Elderly individual maintaining independence, health, and quality of life. Characteristics: active, social, engaged, manageable chronic conditions. Factors: genetics, lifestyle, socioeconomic, healthcare access, social support.

\medterm{Aging}-- Biological process of growing older, involving cumulative changes at molecular, cellular, tissue, organ levels. Hallmarks: genomic instability, telomere attrition, epigenetic alterations, loss of proteostasis, deregulated nutrient sensing, mitochondrial dysfunction, cellular senescence, stem cell exhaustion, altered intercellular communication.

\medterm{Centenarian}-- Person aged 100 or older. ~100,000 worldwide. More women than men (4:1). Often maintain functional independence, cognitive function. Models for healthy aging research.

\medterm{Comprehensive Geriatric Assessment}-- See geriatric assessment.

\medterm{Dentition}-- See dentistry.

\medterm{Elderly}-- Person aged 65 or older. Divisions: young-old (65-74), middle-old (75-84), old-old (85+). Heterogeneous population with varying health status, needs, preferences.

\medterm{Fracture Prevention}-- See osteoporosis.

\medterm{Geriatrix}-- Outdated term for elderly female patient. Prefer "older woman" or "elderly patient."

\medterm{Healthcare Utilization}-- Use of medical services by elderly. Higher rates than other ages. Common: primary care, hospitalization, emergency department, long-term care. Factors: health status, insurance, access, preferences.

\medterm{Incontinence}-- See incontinence.

\medterm{Lifelong Learning}-- Cognitive engagement throughout life. May delay cognitive decline. Activities: reading, puzzles, social interaction, education.

\medterm{Mortality Rate}-- See mortality.

\medterm{Nursing Home}-- See long-term care.

\medterm{Osteoarthritis (Geriatric)}-- Degenerative joint disease, most common arthritis in elderly. Affects weight-bearing joints, hands. Symptoms: pain, stiffness, swelling, decreased range of motion. Treatment: exercise, weight loss, analgesics, joint replacement.

\medterm{Presbycusis}-- Age-related hearing loss. Sensorineural, bilateral, high-frequency. Affects communication, social isolation, cognitive function. Intervention: hearing aids, communication strategies.

\medterm{Pressure Ulcer}-- See pressure injury.

\medterm{Quality of Life}-- See quality of life.

\medterm{Retirement}-- Exit from workforce, typically age 65+. Affects identity, social network, finances, daily structure. Adjustment varies. Planned retirement associated with better outcomes.

\medterm{Screening (Geriatric)}-- Detection of disease in asymptomatic elderly. Balance benefits vs. harms, life expectancy, preferences. Examples: bone density, abdominal aortic aneurysm, colorectal cancer (individualized), depression, falls.

\medterm{Widowhood}-- Loss of spouse. Common in elderly (women more often). Effects: grief, depression, functional decline, mortality increase ("widowhood effect"). Support essential.

\medterm{Fall Prevention}-- See fall.

\medterm{Infection Prevention}-- See infection control.

\medterm{Vaccination}-- See vaccination.

\medterm{Polypharmacy}-- See polypharmacy.

\medterm{Deprescribing}-- Systematic reduction or discontinuation of inappropriate medications. Goals: reduce adverse effects, simplify regimen, improve outcomes.

\medterm{Beers Criteria}-- List of potentially inappropriate medications in elderly. Updated periodically. Guides deprescribing.

\medterm{STOPP/START}-- Screening Tool of Older Person's Prescriptions (STOPP) and Potential Older Person's Prescriptions (START). Identifies inappropriate prescribing.

\medterm{Medication Reconciliation}-- Process of comparing current medications to prescribed medications. Reduces errors at care transitions.

\medterm{Adherence (Geriatric)}-- See general.

\medterm{Cognitive Decline}-- See cognitive impairment.

\medterm{Mild Cognitive Impairment}-- See cognitive impairment.

\medterm{Alzheimer Disease}-- Most common dementia. Progressive neurodegenerative disease. Amyloid plaques, neurofibrillary tangles. Treatment: cholinesterase inhibitors, memantine.

\medterm{Vascular Dementia}-- Second most common dementia. Caused by cerebrovascular disease. Stepwise decline. Treatment: vascular risk factor modification.

\medterm{Lewy Body Dementia}-- Dementia with visual hallucinations, parkinsonism, fluctuating cognition. Alpha-synuclein pathology. Treatment: cholinesterase inhibitors, avoid antipsychotics.

\medterm{Frontotemporal Dementia}-- Dementia affecting frontal and temporal lobes. Behavioral variant (personality changes) or language variant (aphasia). Younger onset (50-60).

\medterm{Delirium}-- See confusion.

\medterm{Dreams (Geriatric)}-- Sleep disturbances common. Nightmares, insomnia. Evaluate for delirium, depression, medications.

\medterm{Sleep (Geriatric)}-- See sleep disorders.

\medterm{Insomnia}-- See sleep disorders.

\medterm{Restless Legs Syndrome}-- Urge to move legs, worse at night. Associated with iron deficiency, renal disease, pregnancy. Treatment: dopamine agonists, iron supplementation.

\medterm{Sleep Apnea}-- See pulmonology.

\medterm{REM Sleep Behavior Disorder}-- Acting out dreams during REM sleep. Associated with Parkinson disease, Lewy body dementia. Treatment: melatonin, clonazepam.

\medterm{Falls}-- See fall.

\medterm{Gait Disturbance}-- See gait.

\medterm{Balance Impairment}-- See balance.

\medterm{Muscle Weakness}-- See sarcopenia.

\medterm{Frailty}-- See frailty.

\medterm{Disability}-- See functional status.

\medterm{Activities of Daily Living}-- See functional status.

\medterm{Instrumental ADL}-- See functional status.

\medterm{Assistive Device}-- See rehabilitation.

\medterm{Home Care}-- Medical care provided at home. Includes nursing, therapy, personal care.

\medterm{Home Health}-- See home care.

\medterm{Skilled Nursing Facility}-- Inpatient facility providing skilled nursing and therapy. For recovery after hospitalization.

\medterm{Assisted Living}-- Residential setting with assistance for ADLs. Less intensive than nursing home.

\medterm{Memory Care}-- Specialized care for dementia patients. Secure environment, structured activities.

\medterm{Hospice}-- See hospice.

\medterm{Palliative Care}-- See palliative care.

\medterm{End-of-Life Care}-- See hospice, palliative care.

\medterm{Advance Care Planning}-- Discussion of future healthcare preferences. Includes advance directives, surrogate designation.

\medterm{Living Will}-- See advance directive.

\medterm{Durable Power of Attorney}-- See advance directive.

\medterm{Healthcare Proxy}-- See advance directive.

\medterm{Do Not Resuscitate}-- DNR: order not to perform CPR. Can be combined with DNI (do not intubate).

\medterm{Comfort Measures Only}-- CMO: focus on symptom relief, not life prolongation.

\medterm{Terminal Illness}-- See prognosis.

\medterm{Pain (Geriatric)}-- See pain management.

\medterm{Dyspnea}-- See pulmonology.

\medterm{Anorexia}-- See nutrition.

\medterm{Malnutrition}-- See clinical nutrition.

\medterm{Dehydration}-- See clinical nutrition.

\medterm{Constipation}-- See gastroenterology.

\medterm{Urinary Retention}-- See urinary retention.

\medterm{Incontinence}-- See incontinence.

\medterm{Osteoporosis}-- See osteoporosis.

\medterm{Osteoarthritis}-- See osteoarthritis.

\medterm{Rheumatoid Arthritis}-- See rheumatology.

\medterm{BPH}-- See BPH.

\medterm{Urinary Tract Infection}-- See CAUTI.

\medterm{Pressure Ulcer}-- See pressure injury.

\medterm{Decubitus Ulcer}-- See pressure injury.

\medterm{Skin Integrity}-- See pressure injury.

\medterm{Wound Care}-- See wound care.

\medterm{Nutrition}-- See nutrition.

\medterm{Swallowing}-- See dysphagia.

\medterm{Dysphagia}-- See dysphagia.

\medterm{Enteral Nutrition}-- See clinical nutrition.

\medterm{Parenteral Nutrition}-- See clinical nutrition.

\medterm{Feeding Tube}-- See enteral nutrition.

\medterm{PEG Tube}-- See medical implants.

\medterm{Quality of Life}-- Perception of well-being including physical, psychological, social, spiritual domains. Measured by QoL instruments.

\medterm{Health-Related QoL}-- See quality of life.

\medterm{Functional Status}-- See functional status.

\medterm{Independence}-- See functional status.

\medterm{Dependence}-- Need for assistance with ADLs/IADLs.

\medterm{Caregiver}-- Person providing care to elderly. Can be family, friend, professional. Burnout common.

\medterm{Caregiver Burden}-- See caregiver stress.

\medterm{Caregiver Stress}-- Emotional, physical, financial strain from caregiving. Screening and support essential.

\medterm{Social Support}-- See social support.

\medterm{Isolation}-- See isolation.

\medterm{Loneliness}-- Subjective feeling of isolation. Associated with depression, cognitive decline, mortality.

\medterm{Social Networks}-- See social support.

\medterm{Community Resources}-- Services available in community: meal delivery, transportation, adult day care, senior centers.

\medterm{Elder Abuse}-- See elder abuse.

\medterm{Neglect}-- See elder abuse.

\medterm{Financial Abuse}-- See elder abuse.

\medterm{Economic Status}-- See social determinants.

\medterm{Poverty}-- See social determinants.

\medterm{Insurance}-- See healthcare access.

\medterm{Medicare}-- See healthcare access.

\medterm{Medicaid}-- See healthcare access.

\medterm{Long-Term Care Insurance}-- See long-term care.

\medterm{Healthcare Access}-- Ability to obtain needed care. Barriers: cost, transportation, availability, literacy.

\medterm{Health Disparities}-- See social determinants.

\medterm{Cultural Competence}-- See communication.

\medterm{Language Barrier}-- See communication.

\medterm{Interpreter}-- See communication.

\medterm{Health Literacy}-- See health literacy.

\medterm{Patient Education}-- See patient education.

\medterm{Self-Management}-- See adherence.

\medterm{Self-Efficacy}-- See self-efficacy.

\medterm{Motivation}-- See behavior change.

\medterm{Behavior Change}-- See behavior change.

\medterm{Smoking Cessation}-- See prevention.

\medterm{Alcohol Use}-- See substance use.

\medterm{Substance Use}-- See substance use.

\medterm{Medication Safety}-- See medication review.

\medterm{Drug Interactions}-- See polypharmacy.

\medterm{Adverse Drug Event}-- See adverse event.

\medterm{Fall Risk}-- See fall.

\medterm{Frailty}-- See frailty.

\medterm{Sarcopenia}-- See sarcopenia.

\medterm{Osteopenia}-- See osteoporosis.

\medterm{Bone Health}-- See osteoporosis.

\medterm{Calcium}-- See biochemistry.

\medterm{Vitamin D}-- See biochemistry.

\medterm{Hormone Therapy}-- See endocrinology.

\medterm{Menopause}-- See obstetrics\_gynecology.

\medterm{Andropause}-- Age-related testosterone decline in men. Symptoms: fatigue, decreased libido, depression. Treatment controversial.

\medterm{Prostate Cancer}-- See oncology.

\medterm{Breast Cancer}-- See oncology.

\medterm{Colorectal Cancer}-- See oncology.

\medterm{Screening}-- See screening.

\medterm{Cancer Screening}-- See screening.

\medterm{Cardiovascular Screening}-- See cardiology.

\medterm{Cognitive Screening}-- See cognitive impairment.

\medterm{Depression Screening}-- See depression.

\medterm{Fall Screening}-- See fall.

\medterm{Bone Density Screening}-- See osteoporosis.

\medterm{Abdominal Aortic Aneurysm Screening}-- See aneurysm.

\medterm{Hearing Screening}-- See presbycusis.

\medterm{Vision Screening}-- See vision impairment.

\medterm{Dental Screening}-- See dentistry.

\medterm{Nutritional Screening}-- See nutritional screening.

\medterm{Functional Screening}-- See functional assessment.

\medterm{Advance Care Planning}-- See advance care planning.

\medterm{Ethics}-- See ethics.

\medterm{Legal Issues}-- See advance directives, elder abuse.

\medterm{Competence}-- See informed consent.

\medterm{Capacity}-- See informed consent.

\medterm{Surrogate Decision-Maker}-- See advance directive.

\medterm{Research Ethics}-- See ethics.

\medterm{Gerontology}-- Study of aging. Biological, psychological, social aspects.

\medterm{Geriatrics}-- Medical care of elderly.

\medterm{Gerontologist}-- Specialist in gerontology.

\medterm{Geriatrician}-- Physician specializing in geriatrics.

\medterm{Aging Population}-- See demographics.

\medterm{Demographics}-- See social determinants.

\medterm{Life Expectancy}-- See life expectancy.

\medterm{Mortality}-- See mortality.

\medterm{Morbidity}-- See morbidity.

\medterm{Healthy Years}-- See healthy years lost.

\medterm{Healthy Aging}-- Process of developing and maintaining functional ability. WHO definition.

\medterm{Successful Aging}-- Model combining low disease, high function, continued engagement.

\medterm{Resilience}-- Ability to recover from stressors. Important for healthy aging.

\medterm{Well-Being}-- See quality of life.

\medterm{Happiness}-- See well-being.

\medterm{Life Satisfaction}-- See well-being.

\medterm{Meaning}-- See well-being.

\medterm{Purpose}-- See well-being.

\medterm{Spirituality}-- See well-being.

\medterm{Religion}-- See spirituality.

\medterm{Existential Distress}-- See psychological.

\medterm{Death}-- End of life. Cultural variations in attitudes.

\medterm{Dying}-- Process approaching death. Physical, psychological, social changes.

\medterm{Bereavement}-- Grief after loss. Normal process but can be complicated.

\medterm{Complicated Grief}-- See bereavement.

\medterm{Funeral}-- See bereavement.

\medterm{Legacy}-- See advance care planning.

\medterm{Wisdom}-- See lifelong learning.
